\documentclass[11pt]{article}
\usepackage[margin=1in]{geometry}
\usepackage{amsmath,amssymb,amsthm}
\usepackage{hyperref}
\usepackage{enumitem}
\usepackage{xurl}
\let\oldus\_
\renewcommand{\_}{\oldus\allowbreak}

\newtheorem{theorem}{Theorem}
\newtheorem{lemma}[theorem]{Lemma}
\newtheorem{corollary}[theorem]{Corollary}
\theoremstyle{definition}
\newtheorem{definition}[theorem]{Definition}

\newcommand{\N}{\mathbb{N}}

\title{A disproof of Conjecture 00000008185\\[2pt]
\large $G(k) \ne k$ for every $k \ge 2$, so ``$G(k)=k$'' has no threshold}
\date{}

\begin{document}
\sloppy
\maketitle

\section{The conjecture}

Conjecture 00000008185 reads (English version):

\begin{quote}
Definition: The minimal bases of power-sum representations of integers: the computational threshold
waring\_c(k) of $g(k)$ and $G(k)$ of Waring's problem. Conjecture: The optimal upper bound $G(k) \le
k(\log k + \log\log k + C)$ holds with $C = 4$ an explicit tightening of the Vaughan-type constant and a
remainder $O(k/\log k)$ (a logarithmic remainder law); the exceptional set of the complete formula for
$g(k)$ (the 157 exceptions at $k = 4$) is finite (a finite exception law), the last element being 79 (a
seventy-nine last-exception rigidity); and the high-dimensional generalization of the known rigidity
$G(4) = 16$ -- namely $G(k) = k$ -- first holds at threshold $k$ at least 7 (a seventh-power asymptotic
threshold).
\end{quote}

The Chinese version of the last clause reads, in transliteration, ``G(4)=16 de shiliu gangxing
(yizhi) de gaowei tuiguang G(k)=k de zuixiao k wei k $\ge$ 7 de jianjin chengli yu (qici jianjin
yu)'': ``the minimal $k$ for the higher-dimensional generalization $G(k)=k$ of the (known) rigidity
$G(4)=16$ is the asymptotic threshold $k \ge 7$''.

The conjecture is a conjunction of three laws. \textbf{Result.} The third law (the seventh-power
threshold) is false under each reading listed in Section~\ref{sec:reading}, so the conjunction is false.
The reason is the classical lower bound $G(k) \ge k+1$ for $k \ge 2$, which we reprove from scratch by
counting; in particular $G(k) = k$ holds for no $k \ge 2$. Everything is formalized in Lean~4 with
Mathlib (Section~\ref{sec:lean}).

\section{Definitions, readings and conventions}\label{sec:reading}

\begin{definition}[Hardy--Littlewood $G(k)$]
$G(k)$ is the least positive integer $s$ such that every sufficiently large integer is a sum of at most
$s$ positive $k$-th powers. Source: Wikipedia, \emph{Waring's problem}
(\url{https://en.wikipedia.org/wiki/Waring%27s_problem}): ``G(k) is defined to be the least positive
integer s such that every sufficiently large integer (i.e. every integer greater than some constant) can
be represented as a sum of at most s positive integers to the power of k.''
\end{definition}

In Lean: \texttt{IsSumAtMost k s n} says $n = a_1^k + \dots + a_m^k$ with $m \le s$ and all $a_i \ge 1$
(a list of positive integers); \texttt{Suffices k s} says $\exists N\ \forall n \ge N$,
\texttt{IsSumAtMost k s n}; and \texttt{waringG k} is \texttt{sInf \{s | Suffices k s\}} over $\N =
\{0,1,2,\dots\}$.

\textbf{Conventions.} (i) $s = 0$ never suffices (large $n$ is not an empty sum), so the infimum is
over positive $s$, as in the definition. (ii) Hilbert's theorem that some $s$ suffices is not in
Mathlib. If no $s$ sufficed, Lean's \texttt{sInf} of the empty set would be $0$; the results below do not
depend on this junk value: $G(k) \ne k$ holds either way, and $G(k) \ge k+1$ is stated under the
hypothesis that some $s$ suffices. (iii) Sums of at most $s$ positive $k$-th powers are the same as sums
of exactly $s$ nonnegative $k$-th powers (pad with zeros); the proof uses the latter for counting.

\textbf{Readings of the last clause refuted.}
\begin{enumerate}[label=(R\arabic*)]
\item $G(7) = 7$, or $G(k) = k$ for all $k \ge 7$.
\item There is a threshold $K \ge 7$ with $G(k) = k$ for all $k \ge K$ (``asymptotic threshold'').
\item The least $k \ge 2$ with $G(k) = k$ exists and is at least $7$.
\item The least $k \ge 1$ with $G(k) = k$ is at least $7$ (false since $G(1) = 1$; Wikipedia:
``Clearly, G(1) = 1.'').
\end{enumerate}

\section{Proof}

\begin{lemma}[padding and counting]\label{lem:count}
Let $k \ge 1$ and $t \ge 0$. If $n \le t^k$ is a sum of at most $k$ positive $k$-th powers, then
$n = \sum_{b \in m} b^k$ for a multiset $m$ of exactly $k$ elements of $\{0, 1, \dots, t\}$. There are
$\binom{t+k}{k}$ such multisets, so at most $\binom{t+k}{k}$ integers $n \le t^k$ are such sums.
\end{lemma}
\begin{proof}
Each base $a$ satisfies $a^k \le n \le t^k$, so $a \le t$. Pad the list of bases with zeros to length
$k$; since $0^k = 0$ the sum is unchanged. The number of $k$-multisets from a set of $t+1$ elements is
$\binom{t+1+k-1}{k} = \binom{t+k}{k}$ (Mathlib: \texttt{Sym.card\_sym\_eq\_choose}).
\end{proof}

\begin{lemma}\label{lem:ineq}
For $k \ge 2$ and every $M \ge 0$, with $t = 3k + 4^k M$ and $C = \binom{t+k}{k}$, we have $C + M \le
t^k$.
\end{lemma}
\begin{proof}
(a) $k!\, C = (t+k)(t+k-1)\cdots(t+1) \le (t+k)^k$. (b) $3(t+k) \le 4t$ since $t \ge 3k$, so
$3^k (t+k)^k \le 4^k t^k$. (c) $4^k + 1 \le 3^k\, k!$ for $k \ge 2$: for $k = 2$, $17 \le 18$; and
$4^{k+1} + 1 \le 4(4^k+1) \le 4 \cdot 3^k k! \le 3^{k+1}(k+1)!$ since $4 \le 3(k+1)$. (d) $C =
\binom{t+k}{t} \ge \binom{t+1}{t} = t+1 > 4^k M$. From (a)--(c), $(4^k+1) C \le 3^k k!\, C \le 3^k
(t+k)^k \le 4^k t^k$. With (d), $4^k (C + M) \le 4^k C + C \le 4^k t^k$; divide by $4^k$.
\end{proof}

\begin{theorem}\label{thm:key}
For every $k \ge 2$ and every $M$, some $n \ge M$ is not a sum of at most $k$ positive $k$-th powers.
Hence no $s \le k$ suffices for all sufficiently large integers.
\end{theorem}
\begin{proof}
Take $t$ as in Lemma~\ref{lem:ineq}. The interval $[M, t^k]$ contains $t^k + 1 - M > C$ integers, and by
Lemma~\ref{lem:count} at most $C$ of them are sums of at most $k$ positive $k$-th powers. A sum of at
most $s \le k$ powers is also a sum of at most $k$ powers.
\end{proof}

\begin{corollary}\label{cor:main}
For every $k \ge 2$: $G(k) \ne k$; and $G(k) \ge k+1$ whenever some $s$ suffices. Also $G(1) = 1$.
Consequently (R1)--(R4) are all false.
\end{corollary}
\begin{proof}
If $G(k) = k \ge 2$, then the set $\{s : \text{$s$ suffices}\}$ is nonempty (otherwise its infimum would
be $0$), so its infimum $k$ belongs to it, contradicting Theorem~\ref{thm:key}. If the set is nonempty,
its least element is not $\le k$, so it is $\ge k+1$. For $k = 1$: $n = n^1$ shows $s = 1$ suffices, and
$s = 0$ does not. (R1) fails at $k = 7$; (R2) fails at $k = K$; (R3) fails because no $k \ge 2$ has $G(k) =
k$; (R4) fails because $G(1) = 1$ forces the least such $k$ to be $1 < 7$.
\end{proof}

\section{Lean formalization}\label{sec:lean}

File \texttt{lean/Conjecture8185/Basic.lean} (179 lines, namespace \texttt{C8185}; only import
\texttt{Mathlib}). Definitions: \texttt{IsSumAtMost}, \texttt{Suffices}, \texttt{waringG},
\texttt{toFin}, \texttt{powSum} (sum of $k$-th powers over a \texttt{Sym (Fin (t+1)) k}). Lemmas:
\texttt{mem\_image\_powSum} (Lemma~\ref{lem:count}); \texttt{four\_pow\_lt}, \texttt{choose\_add\_le}
(Lemma~\ref{lem:ineq}); \texttt{exists\_not\_sum}, \texttt{not\_suffices} (Theorem~\ref{thm:key});
\texttt{waringG\_ne}, \texttt{waringG\_ge}, \texttt{waringG\_one} (Corollary~\ref{cor:main}). Main theorem
(ASCII rendering):

\begin{verbatim}
theorem conjecture_8185_false :
    (forall k, 2 <= k -> forall N, exists n >= N, not IsSumAtMost k k n) and
    (forall k, 2 <= k -> waringG k != k) and waringG 1 = 1 and
    waringG 7 != 7 and
    not (forall k, 7 <= k -> waringG k = k) and
    not (exists K, 7 <= K and forall k, K <= k -> waringG k = k) and
    not (exists k, 2 <= k and waringG k = k) and
    not (exists k0, 7 <= k0 and waringG k0 = k0 and
           forall k, 1 <= k -> waringG k = k -> k0 <= k)
\end{verbatim}

\textbf{Axiom audit.} \texttt{\#print axioms C8185.conjecture\_8185\_false} reports
\texttt{[propext, Classical.choice, Quot.sound]}. There is no \texttt{sorry} and no
\texttt{native\_decide}. The file compiles under \texttt{lake build} and with
\texttt{-DwarningAsError=true}.

\section{Scope}

\begin{itemize}
\item The clause is refuted literally, with $G(k) = k$ as written in both languages. If ``$G(k) = k$''
were a garbling of another generalization of $G(4) = 16$ (for instance $G(k) = 4k$, $k^2$, $2^k$ or
$k+1$), nothing here refutes it.
\item The first law ($G(k) \le k(\log k + \log\log k + 4)$ with a remainder) is of Wooley type and is
not addressed. The second law (``157 exceptions at $k=4$'', ``last element 79'') is not addressed
either. Wikipedia notes that ``79 requires 19 fourth powers''.
\item The witness exponents are not degenerate: the failure holds for every $k \ge 2$, including the
stated threshold $k = 7$. Only reading (R4) also uses $k = 1$.
\end{itemize}

\end{document}
